% LSPV-2 REPORT::CUBEROOTSEL1::2026-10-14T00:28:52+09:00
% !TEX program = xelatex
\documentclass[11pt]{article}
\usepackage[T1]{fontenc}
\usepackage{lmodern,amsmath,amssymb,amsthm,booktabs,tabularx,xurl,hyperref}
\usepackage[margin=0.84in]{geometry}
\newtheorem{theorem}{Theorem}
\newtheorem{corollary}[theorem]{Corollary}
\newcommand{\F}{\mathbb F_2}
\newcommand{\EXACT}{\textbf{PROVED/EXACT}}
\newcommand{\KNOWN}{\textbf{KNOWN THEOREM}}
\newcommand{\COMP}{\textbf{COMPUTATIONAL}}
\newcommand{\COND}{\textbf{CONDITIONAL}}
\newcommand{\CAND}{\textbf{CANDIDATE}}
\newcommand{\OPEN}{\textbf{OPEN}}
\newcommand{\RET}{\textbf{FALSIFIED/RETIRED}}
\hypersetup{colorlinks=true,urlcolor=blue,linkcolor=blue}
\setlength{\emergencystretch}{3em}
\title{CUBEROOTSEL1: A Target-Only Exponent Selector for\\
Non-Surjective Even-Degree Cube Residuals}
\author{UNIVERSAL\_CORE\_TESTS / LSPV-2}
\date{14 October 2026, 00:28 KST}

\begin{document}
\maketitle

\begin{center}
\fbox{\parbox{0.95\linewidth}{
\textbf{Executive status: P versus NP remains OPEN.}
PLANTEDCUBEFW1 correctly gives a source-witness section for every even
degree.  On the infinite subtype $d\not\equiv0\pmod6$, however, the planted
preimage is unnecessary.  If $m=(2^d-1)/3$ and
$e=3^{-1}\bmod m$, then $b^e$ is a cube root of every promised nonzero image
target $b$.  Exponentiation followed by one cube check also recognizes the
image.  The selector survives explicit invertible affine input/output mixing
and complete Tseitin forest lifting.  Thus ``remove the planted section'' is
not a sufficient repair for this cube family.  The symbolic theorem is
\CAND{} pending independent review; the finite audit is \COMP{}.}}
\end{center}

\section{Claim card}
\begin{tabularx}{\linewidth}{p{0.27\linewidth}X}
\toprule
Claim ID & CUBEROOTSEL1 \\
Status before run & PLANTEDCUBEFW1 theorem \EXACT{} after independent review; its finite audit \COMP{} \\
Claim attacked & Removing the supplied/planted source section may leave a hard image-promised even-cube preimage problem \\
Exact quantifiers & Every even $d\ge4$ with $6\nmid d$; every declared irreducible polynomial basis of $K=\mathbb F_{2^d}$; every $b\in\operatorname{im}(z\mapsto z^3)$ \\
Machine/model & Uniform deterministic random-access bit machine, full typed-record access, polynomial workspace \\
Resource & Parsing, exponent construction, field operations, image check, affine decoding, forest reconstruction, verification, and output \\
SAME\_REALIZATION & The same target $b$, public field/affine record, selected root $z$, charge vector, and reconstructed Tseitin edge witness \\
Falsifier & An image target for which the exponent output fails to cube to $b$, or a returned forest witness violating incidence/residual constraints \\
\bottomrule
\end{tabularx}

The selected obstacles were (i) whether a promised image point can be inverted
without the source witness or hidden preimage and (ii) whether this selector
survives the actual cotree/forest interface.  The weakest useful advance was
an explicit right inverse on one infinite, polynomially scalable subtype.

\section{Primary literature and transfer classification}
Adleman, Manders, and Miller treat root extraction in finite fields as a
computational problem in its own right \cite{AMM}.  Barreto and Voloch give
efficient root algorithms under arithmetic conditions on the field and root
degree \cite{BV}.  Their directly transferable lesson is algorithmic: a
many-to-one power map can still admit efficient promised-image root
extraction.  Neither paper gives an unrestricted SAT lower bound.

The theorem below needs no imported complexity claim; it follows by a cyclic
group calculation.  Finite-field power extraction is therefore an
\textbf{EXACT OBJECT/THEOREM TRANSFER}.  Treating non-surjectivity, collision
entropy, or lack of a total inverse as runtime hardness is a
\textbf{FAILED TRANSFER}.  Schur powers, quadratic-relation spaces, syzygies,
and LWE-style worst-case/average-case architectures remain
\textbf{ANALOGY/CANDIDATE}; no such bridge is used here.

\section{Target-only cube selector}
Let $d\ge4$ be even, $K=\mathbb F_{2^d}$, and
\[
 P_d(z)=z^3,\qquad m=\frac{2^d-1}{3}.                   \tag{1}
\]
The nonzero image $P_d(K^*)$ is the unique subgroup of $K^*$ of order $m$.

\begin{theorem}[promised-image exponent selector]
If $6\nmid d$, then $\gcd(3,m)=1$.  Let
\[
 e\equiv3^{-1}\pmod m.                                 \tag{2}
\]
Define $R_d(0)=0$ and $R_d(b)=b^e$ for $b\ne0$.  For every
$b\in\operatorname{im}P_d$,
\[
 P_d(R_d(b))=b.                                         \tag{3}
\]
Moreover, computing $z=R_d(b)$ and accepting exactly when $z^3=b$ is a
uniform deterministic polynomial-time image recognizer and selector.
\end{theorem}
\begin{proof}
The powers of $2$ modulo $9$ have period six.  For even $d$, if $6\nmid d$
then $2^d\equiv4$ or $7\pmod9$, so $m$ is not divisible by three and (2)
exists.  For nonzero $b\in\operatorname{im}P_d$, we have $b^m=1$.  Write
$3e=1+tm$.  Then
\[
 (b^e)^3=b^{3e}=b^{1+tm}=b.
\]
Zero is immediate.  Binary exponentiation uses an $O(d)$-bit exponent and
$O(d)$ field multiplications; polynomial-basis field arithmetic and the final
cube check are polynomial in the typed encoded length.  If the check fails,
$b$ was not accepted as an image point.
\end{proof}

This is a right inverse only on the image.  It does not contradict the
three-to-one nonzero fibers or create a total inverse $K\to K$.

\paragraph{Exact boundary.}
If $6\mid d$, then $2^d\equiv1\pmod9$ and $3\mid m$, so no exponent $e$ can
satisfy $3e\equiv1\pmod m$.  This retires only the monomial exponent selector
on those dimensions.  It is not a lower bound against other deterministic
root algorithms.

\section{Affine conjugates and the Tseitin compiler subtype}
Let $A,B\in\operatorname{GL}_d(\F)$ and $p,q\in\F^d$ be public typed data.
For
\[
 Q(z)=B\bigl((Az+p)^3\bigr)+q,                          \tag{4}
\]
an image target $w$ has the selected preimage
\[
 z=A^{-1}\bigl(R_d(B^{-1}(w+q))+p\bigr).               \tag{5}
\]
Thus coordinate permutation, complementation/phase, and dense invertible
linear mixing do not remove the selector when their acquisition is charged.

For the $K_{s,s}$ cotree used in PLANTEDCUBEFW1, $d=(s-1)^2$.  Choose the
least odd $s$ with
\[
 (s-1)^2\ge\max\{4,n\},\qquad 3\nmid(s-1).             \tag{6}
\]
Then $d\equiv4\pmod6$, $s\le\lceil\sqrt n\rceil+4$, and $d=O(n)$.
Consequently (6) gives an infinite polynomial-size typed compiler subtype.
After any source witness is found, evaluate the declared residual target,
apply (5), assign the selected value to the cotree, and solve each spanning
tree leaf-to-root.  The residual preimage need not be included in the record
or recovered from the original planted section.

\begin{corollary}[charged residual neutrality]
On the subtype (6), every promised satisfiable residual target has a uniform
target-only polynomial selector, followed by a polynomial forest lift.  The
residual layer therefore contributes only charged polynomial work beyond
finding the embedded source witness.
\end{corollary}

The result does not solve the source SAT instance.  It sharpens the next
obligation: a candidate must defeat both source-witness sections and all
target-only promised-image selectors.

\section{Bounded reproducible audit}
The dependency-free program
\path{experiments/review/even_cube_target_selector_audit.py} has SHA-256
\path{12a2670a6ad54888f7a2b505aac07a5528bbb492a419948e686fe6361cb6fbe3}.
One deterministic invocation completed in 1.143 seconds and wrote
\path{results/raw/review/2026-10-14_0028_even_cube_target_selector.json}.

\begin{center}
\begin{tabular}{lr}
\toprule
Check & Count \\
\midrule
Assertions & 27,810 \\
Image targets selected & 5,896 \\
Off-image targets rejected & 11,784 \\
Affine-conjugate targets selected & 1,296 \\
Recovered roots differing from erased planted root & 862 \\
Complete Tseitin witnesses & 540 \\
Source YES / NO controls & 12 / 4 \\
Failures & 0 \\
\bottomrule
\end{tabular}
\end{center}

Degrees $d=4,8,10,14$ were exhaustive.  Their selector exponents were
$2,57,114,3641$, and their image sizes were $6,86,342,5462$, with exactly one
zero fiber of size one and every nonzero image fiber of size three.  The
actual square cotree dimensions $d=4,16$ supplied 528 preimage-erased full
witness reconstructions; twelve additional source-YES controls gave 540.
The audit also covered public dense affine mixing, phases, edge renaming,
disconnected graphs with an isolate, odd component charge, reducible modulus,
odd field degree, and the $d=6,12,18$ no-power-exponent boundary.  No
invocation required chunking or approached thirty seconds.  Wall time is not
used as asymptotic evidence.

\section{Hostile controls, barriers, and PQ ledger}
The selector permits full source and target access, polynomial workspace,
Gaussian/affine preprocessing, compact field arithmetic, and input-assisted
reconstruction.  Horn, dual-Horn, Krom, affine, bounded-treewidth,
independent-product, padding, renaming, complementation, DEAD, COPY, DELAY,
INTERP, and RECOMP controls cannot make the explicit selector expensive; a
solver may simply skip irrelevant wrappers.  A maximal semantic or algebraic
oracle only strengthens this upper-bound escape.

The result relativizes and algebrizes and supplies no nonrelativizing hinge.
Natural Proofs is not invoked.  DNNF, PC/PCR, Resolution, pebbling,
communication, and one-tape bounds remain model-specific.  No ETH,
cryptographic, random-oracle, or average-case MQ assumption is used.

The PQ classification is: the cyclic image subgroup and power selector are
an \textbf{EXACT OBJECT/THEOREM TRANSFER}; collision entropy is algorithmically
exploitable rather than hard here; Schur-square, relation-rank, syzygy, and
special-soundness analogies are \textbf{FAILED TRANSFERS} to unrestricted
runtime absent a new same-realization theorem.

\section{CORPUS REVIEWED}
The run surveyed the current GitHub repository tree and read root
\path{main.tex}, AGENTS.md, CODEX.md, the complete PNP research skill and its
experiment/protocol references, all seven current task specifications,
AFFINE-RESIDUAL-REDIRECT1, EARLY22-INTEGRATION1, SINGLE22-SEPARATION1, and the
AFFINE-RESIDUAL-RANK1 exploration.  It read the newest verified MAIN
PROMISEGATE1 and the relevant CUBEPERMFW1/FIELDTRACEFW1 snapshots.  Material
specialist sources included PLANTEDCUBEFW1 and its code/results,
CHECKABLEPERM1, ORBITDECODERFW1, HIERLIFT-TEST1, FIELDTRACELIFT1, PERMEXT1,
POLARKERNEL1, QUADSHEAR1, and BIJECTIVEQFW1.  HEC/CRS/RCCN/Search-G task
specifications were inspected as transfer guards; their restricted online
lower bounds do not apply to this explicit upper-bound selector.

Predecessor recovery was not required: the newest LSPV-1 report
\path{research/agents/lspv-1/2026-10-14_0007_LSPV-1_PLANTEDCUBEFW1.tex}
is substantive, has the required header/sentinel, and was read back from
GitHub at blob
\path{1171db21fc748946fc3eca4517f7911fa9290b12}.

\section{DUPLICATION / CONTRADICTION AUDIT}
CUBEPERMFW1 handles odd $d$, where cubing is a total permutation.
PLANTEDCUBEFW1 handles every even $d$ using a source-computable section.
Repository search found no prior project report giving the image-subgroup
exponent (2) for the even, non-surjective subtype.  The new result is not a
contradiction: it is a second, source-independent selector on
$d\not\equiv0\pmod6$.

No conflict arises with root-extraction literature, general MQ difficulty,
SAT NP-completeness, or the EARLY22/SINGLE22 exact-$(2,2)$ compiler and
restricted proof-system theorems.  Expanded-monomial proof size remains
distinct from arbitrary deterministic witness search.

\section{WHAT CHANGED RELATIVE TO ROOT MAIN}
Root MAIN correctly states that the planted even-cube family has a polynomial
source-witness section.  This run strengthens the firewall on an infinite
polynomially scalable subtype: the residual root can be selected from the
evaluated target alone.  The next target must therefore exclude both planted
sections and direct image-subgroup power selectors.

\section{ADOPTED / REJECTED / RETIRED PRIOR COMPONENTS}
\begin{itemize}
\item \textbf{Adopted:} the even-cube image/fiber theorem, strict typed field
and graph recognition, cotree/forest lifting, and full cost accounting.
\item \textbf{Preserved:} exact-$(2,2)$ compilation and PC/PCR/Resolution
results only within their verified models; the conditional growing-factor
lemma remains uninstantiated.
\item \textbf{Rejected:} hidden nonresidue advice, uncharged field
decomposition, wall-time extrapolation, and failure of a single exponent on
$6\mid d$ as evidence of hardness.
\item \textbf{Retired:} removal or erasure of the planted preimage as a
sufficient hardness repair for the subtype $6\nmid d$.
\end{itemize}

\section{STALE OR OVERSTATED CLAIMS FOUND}
The phrase ``no preimage is included or planted'' is not enough to specify a
hard section problem.  For a promised image point, a right inverse may be
computed directly from the target.  Conversely, the absence of the exponent
(2) when $6\mid d$ does not show that those instances are hard; it closes only
one selector formula.

\section{Status and handoff}
\begin{center}
\begin{tabularx}{0.98\linewidth}{p{0.39\linewidth}p{0.20\linewidth}X}
\toprule
Claim & Status & Scope \\
\midrule
Target-only image exponent selector & \CAND & Complete elementary proof; independent review pending \\
Finite field/affine/Tseitin audit & \COMP & 27,810 bounded assertions, zero failures \\
Planted-section erasure as cube repair & \RET & Selector ignores the erased preimage for $6\nmid d$ \\
$6\mid d$ general deterministic cube-root cost & \OPEN & No power exponent; no lower bound follows \\
Unrestricted all-program lower bound & \OPEN & No model-matched transfer \\
P versus NP separation & \OPEN & A load-bearing implication remains open \\
\bottomrule
\end{tabularx}
\end{center}

\textbf{Single highest-value next target.}
Work first on the remaining $d\equiv0\pmod6$ cube subtype, but attack it as an
algorithm problem: implement deterministic cubic-root methods, subfield/tower
decomposition, and fixed typed nonresidue search.  If that subtype also has a
uniform selector, retire the entire even-cube route.  If it does not, absence
of known algorithms must still not be promoted to a lower bound.

\begin{thebibliography}{9}
\bibitem{AMM}
L. M. Adleman, K. L. Manders, and G. L. Miller,
``On Taking Roots in Finite Fields,''
\emph{18th Annual Symposium on Foundations of Computer Science}, 1977,
pp. 175--178. doi:10.1109/SFCS.1977.18.

\bibitem{BV}
P. S. L. M. Barreto and J. F. Voloch,
``Efficient Computation of Roots in Finite Fields,''
\emph{Designs, Codes and Cryptography} 39(2), 2006, pp. 275--280.
doi:10.1007/s10623-005-4017-5.
\end{thebibliography}

\bigskip
\noindent P versus NP remains open.

\end{document}
% END_OF_VERIFIED_REPORT::LSPV-2::CUBEROOTSEL1
